\chapter{Dasar Teori}

\section{Konsep Dasar Otomatisasi Perangkat Lunak}

\subsection{\textit{Continuous Integration} dan \textit{Continuous Deployment} (CI/CD)}
\textit{Continuous Integration} (CI) merupakan praktik otomatisasi dalam rekayasa perangkat lunak di mana setiap perubahan kode digabungkan dan diuji secara berkala. Sedangkan \textit{Continuous Deployment} (CD) adalah kelanjutan dari CI yang secara otomatis meluncurkan (\textit{deploy}) kode yang telah lolos uji ke lingkungan produksi. Kombinasi CI/CD mempercepat rilis perangkat lunak sekaligus meminimalisasi \textit{human error} saat peluncuran [\ref{fowler}].

\begin{figure}[H]
    \centering
    \includegraphics[width=0.9\textwidth]{cicd-diagram.png}
    \caption{Ilustrasi alur siklus \textit{Continuous Integration} dan \textit{Continuous Deployment}}
\end{figure}

\subsection{Penerapan CI/CD pada Framework Laravel}
Penerapan otomatisasi pada \textit{framework} Laravel (PHP) umumnya melibatkan instalasi dependensi menggunakan \texttt{Composer}, penyiapan \textit{environment} (\texttt{.env}), serta eksekusi pengujian menggunakan PHPUnit atau perintah bawaan \texttt{php artisan test}. Pada tahap \textit{deployment}, skrip CI/CD juga bertugas menjalankan proses optimasi (\textit{caching}, \textit{route}, \textit{view}) agar performa aplikasi Laravel tetap maksimal saat berada di \textit{server} produksi [\ref{github1}].

\section{Infrastruktur dan Keamanan \textit{Deployment}}

\subsection{GitHub Actions untuk Otomatisasi \textit{Deployment}}
GitHub Actions adalah platform CI/CD yang terintegrasi langsung di repositori GitHub. \textit{Workflow} dibangun menggunakan bahasa YAML (\texttt{.yml}) dan dieksekusi oleh \textit{runner}. Pada proyek Laravel, \textit{workflow} dapat dipadukan dengan skrip \textit{bash} kustom (seperti \texttt{deploy.sh}) untuk mengeksekusi serangkaian perintah \textit{server} secara berurutan.

\begin{figure}[H]
    \centering
    \includegraphics[width=0.9\textwidth]{github-actions.png}
    \caption{Konsep arsitektur \textit{workflow} dan eksekusi pada GitHub Actions}
\end{figure}

\subsection{\textit{Environment Protection Rules}}
Mendeploy aplikasi secara langsung ke \textit{Production} memiliki risiko yang tinggi. Oleh karena itu, GitHub menyediakan fitur \textit{Environment Protection Rules} yang dapat menahan aliran rilis aplikasi hingga mendapatkan persetujuan manual (ACC) dari \textit{reviewer} yang berwenang [\ref{github2}].

\chapter{Hasil dan Pembahasan}

\section{Persiapan Lingkungan dan Konfigurasi}

\subsection{Setup Repositori Sinkronisasi}
Pada tahap awal, repositori organisasi mata kuliah (KEPL2026) sedang mengalami \textit{flagging} sehingga GitHub Actions tidak aktif. Untuk menyiasatinya, dibuat repositori mandiri baru (\texttt{AcademyCity-RND/}\allowbreak\texttt{evolusi-pl-24-542159-SV-25009-v2}) yang dikonfigurasi agar tersinkronisasi (\textit{multi-remote}) dengan repositori asal. Kondisi awal \textit{branch} \texttt{main} dari sisa pertemuan sebelumnya terlihat pada gambar berikut.

\begin{figure}[H]
    \centering
    \includegraphics[width=0.9\textwidth]{6.png}
    \caption{Tampilan awal branch main kondisi awal (bekas pertemuan 1)}
\end{figure}

\subsection{Konfigurasi Script CI/CD Laravel}
Pusat eksekusi \textit{Continuous Deployment} ini dijalankan melalui skrip \texttt{deploy.sh} dan aturan \textit{trigger} dari \texttt{ci.yml}. Skrip \texttt{deploy.sh} bertugas menjalankan perintah operasional seperti memanajemen repositori di \textit{server} tujuan.

\begin{figure}[H]
    \centering
    \includegraphics[width=0.9\textwidth]{1.png}
    \caption{Isi dari script deploy.sh}
\end{figure}

File \texttt{ci.yml} diatur ke dalam 4 tahapan berurutan, yaitu: \textit{Build}, \textit{Test}, \textit{Staging}, dan \textit{Production}. Tahap \textit{Build} dan \textit{Test} dikonfigurasi khusus untuk menyetel \textit{environment} agar pengujian aplikasi berjalan dengan benar.

\begin{figure}[H]
    \centering
    \includegraphics[width=0.9\textwidth]{2.png}
    \caption{Konfigurasi tahap Build}
    
    \vspace{0.5cm}
    
    \includegraphics[width=0.9\textwidth]{3.png}
    \caption{Konfigurasi tahap Test}
\end{figure}

\begin{figure}[H]
    \centering
    \includegraphics[width=0.9\textwidth]{4.png}
    \caption{Lanjutan konfigurasi tahap Test}
    
    \vspace{0.5cm}
    
    \includegraphics[width=0.9\textwidth]{5.png}
    \caption{Konfigurasi tahap Staging dan Production}
\end{figure}

\section{Eksekusi dan Pengujian \textit{Workflow}}

\subsection{Uji Coba CI pada Branch \texttt{feat}}
Saat fitur baru dikembangkan pada \textit{branch} \texttt{feat}, tahap \textit{Test} sengaja disalahkan. Sistem CI langsung menangkap kegagalan tersebut dan memberikan peringatan (warna merah). Setelah skrip \textit{test} diperbaiki, status GitHub Actions kembali hijau.

\begin{figure}[H]
    \centering
    \includegraphics[width=0.9\textwidth]{7.png}
    \caption{Test sengaja digagalkan pada branch feat}
    
    \vspace{0.5cm}
    
    \includegraphics[width=0.9\textwidth]{8.png}
    \caption{Test berhasil diperbaiki pada branch feat}
\end{figure}

Sesuai aturan di \texttt{ci.yml}, peluncuran (\textit{deployment}) ke \textit{Staging} dan \textit{Production} dibatalkan (\textit{skip}) karena pengujian masih dilakukan di lingkup \textit{branch} \texttt{feat}.

\subsection{Uji Coba CI pada Branch \texttt{dev}}
Selanjutnya, kode diintegrasikan ke dalam \textit{branch} \texttt{dev}. Eksekusi berjalan sukses hingga tahap \textit{Staging}. Tahap \textit{Production} masih tertahan (\textit{skip}) untuk menjaga kestabilan rilis.

\begin{figure}[H]
    \centering
    \includegraphics[width=0.9\textwidth]{9.png}
    \caption{Hasil eksekusi sukses hingga tahap Staging pada branch dev}
\end{figure}

\clearpage
\section{Pengamanan Rilis dan Analisis}

\subsection{Pengaturan \textit{Environment Production}}
Untuk menjamin keamanan \textit{Continuous Deployment}, \textit{environment} khusus bernama \texttt{production} ditambahkan langsung pada konfigurasi GitHub.

\begin{figure}[H]
    \centering
    \includegraphics[width=0.9\textwidth]{18.png}
    \caption{Membuat Environment baru di menu Settings}
    
    \vspace{0.5cm}
    
    \includegraphics[width=0.9\textwidth]{19.png}
    \caption{Penamaan Environment production}
\end{figure}

Di dalam pengaturan ini, persetujuan manual (\textit{Required reviewers}) diaktifkan dan menugaskan sebuah akun khusus untuk memberikan persetujuan sebelum \textit{deployment} dieksekusi.

\begin{figure}[H]
    \centering
    \includegraphics[width=0.9\textwidth]{20.png}
    \caption{Konfigurasi Required reviewers dan akun GitHub peninjau}
    
    \vspace{0.5cm}
    
    \includegraphics[width=0.9\textwidth]{21.png}
    \caption{Penyimpanan aturan dan notifikasi keberhasilan}
\end{figure}

\subsection{Eksekusi CI pada Branch \texttt{main}}
Puncak alur kerja ini terjadi saat penggabungan ke \textit{branch} \texttt{main}. Tahapan \textit{Build}, \textit{Test}, dan \textit{Staging} berjalan otomatis, namun \textit{deployment} \textit{Production} tertahan (\textit{pending}).

\begin{figure}[H]
    \centering
    \includegraphics[width=0.9\textwidth]{10.png}
    \caption{Tahap Production tertahan menunggu review}
\end{figure}

\begin{figure}[H]
    \centering
    \includegraphics[width=0.9\textwidth]{11.png}
    \caption{Proses pemberian persetujuan oleh Reviewer}
    
    \vspace{0.5cm}
    
    \includegraphics[width=0.9\textwidth]{12.png}
    \caption{Keempat tahap berhasil dieksekusi pasca persetujuan}
\end{figure}

\subsection{Analisis Log Eksekusi}
Semua tindakan \textit{runner} dicatat dalam log konsol, membuktikan validitas skrip dan \textit{pipeline} yang telah disusun.

\begin{figure}[H]
    \centering
    \includegraphics[width=0.9\textwidth]{13.png}
    \caption{Log output Test yang sengaja digagalkan}
\end{figure}

\begin{figure}[H]
    \centering
    \includegraphics[width=0.9\textwidth]{14.png}
    \caption{Log eksekusi tahap Build sukses}
    
    \vspace{0.5cm}
    
    \includegraphics[width=0.9\textwidth]{15.png}
    \caption{Log eksekusi tahap Test sukses}
\end{figure}

\begin{figure}[H]
    \centering
    \includegraphics[width=0.9\textwidth]{16.png}
    \caption{Log eksekusi tahap Staging sukses}
    
    \vspace{0.5cm}
    
    \includegraphics[width=0.9\textwidth]{17.png}
    \caption{Log eksekusi tahap Production sukses}
\end{figure}

\subsection{Riwayat \textit{Pull Request}}
Alur CI/CD yang komprehensif ini tercermin utuh dalam riwayat penggabungan kode (\textit{merge}) dari awal fase pengembangan hingga siap meluncur ke \textit{server} produksi.

\begin{figure}[H]
    \centering
    \includegraphics[width=0.9\textwidth]{22.png}
    \caption{Riwayat Pull Request dari feat ke dev}
    
    \vspace{0.5cm}
    
    \includegraphics[width=0.9\textwidth]{23.png}
    \caption{Riwayat Pull Request dari dev ke main}
\end{figure}

\chapter{Kesimpulan}
Praktikum Konstruksi dan Evolusi Perangkat Lunak (KEPL) Pertemuan 3 telah sukses mengimplementasikan alur \textit{Continuous Deployment} (CD) pada proyek berbasis Laravel dengan menggunakan GitHub Actions. Kendala pada repositori pusat berhasil diatasi dengan teknik sinkronisasi repositori ganda, memastikan berjalannya otomasi tanpa hambatan. \textit{Workflow} CI/CD terbagi menjadi 4 fase krusial: \textit{Build, Test, Staging,} dan \textit{Production}, di mana setiap fase dieksekusi secara kondisional berdasarkan pergerakan \textit{branch} (\texttt{feat}, \texttt{dev}, dan \texttt{main}). Sistem proteksi berhasil diterapkan secara kuat dengan fitur \textit{Environment Protection Rules}, sehingga \textit{deployment} Laravel ke lingkungan \textit{Production} wajib melewati proses persetujuan (ACC) manual dari akun peninjau. Hal ini mendemonstrasikan praktik terbaik dari integrasi berkelanjutan yang aman dan sistematis.

\chapter{Daftar Pustaka}

\begin{enumerate}
    % Fundamental konsep Continuous Integration/Continuous Deployment
    \item \label{fowler} M. Fowler. \textit{Continuous Integration}. MartinFowler.com, 2006. [Online]. Tersedia: \url{https://martinfowler.com/articles/continuousIntegration.html}.
    
    % Dokumentasi resmi tentang struktur GitHub Actions
    \item \label{github1} GitHub Docs. \textit{Understanding GitHub Actions}. [Online]. Tersedia: \url{https://docs.github.com/}.
    
    % Dokumentasi spesifik mengenai Environment Protection Rules
    \item \label{github2} GitHub Docs. \textit{Using environments for deployment}. [Online]. Tersedia: \url{https://docs.github.com/}.
    
    % Repositori Source Code tugas praktikum
    \item \label{sourcecode} Source Code Praktikum:
    \begin{itemize}
        \item Repo Matkul KEPL2026:\\
        \url{https://github.com/KEPL2026/evolusi-pl-24-542159-SV-25009-v1}
        
        \item Repo Pribadi:\\
        \url{https://github.com/AcademyCity-RND/evolusi-pl-24-542159-SV-25009-v2}
    \end{itemize}
\end{enumerate}
